% 2023 全国甲卷文科 6：程序框图
% 源图流程：输入 n=3,A=1,B=2,k=1；当 k<=n 时依次更新 A、B、k 并折返判断；
% 条件不成立时输出 B。源图所有边均为实线。
\begin{tikzpicture}[font=\footnotesize, >=Stealth, node distance=2.2mm, line width=0.4pt,
  box/.style={draw, minimum height=1.5em, minimum width=5.8em, inner xsep=6pt},
  term/.style={draw, rounded corners=6pt, minimum height=1.5em, inner xsep=8pt},
  io/.style={draw, trapezium, trapezium left angle=70, trapezium right angle=110,
             minimum height=1.5em, inner xsep=3pt},
  dec/.style={draw, diamond, aspect=2.6, inner sep=1.5pt}]
  \node[term] (start) {开始};
  \node[io, below=of start] (input) {输入 $n=3,A=1,B=2,k=1$};
  \coordinate[below=1mm of input] (merge);
  \node[dec, below=1mm of merge] (cond) {$k\leqslant n$};
  \node[box, below=of cond] (stepa) {$A=A+B$};
  \node[box, below=of stepa] (stepb) {$B=A+B$};
  \node[box, below=of stepb] (stepk) {$k=k+1$};
  \node[io, below=2.5mm of stepk] (output) {输出 $B$};
  \node[term, below=of output] (end) {结束};

  % 主流程与“是”分支
  \draw[->] (start) -- (input);
  \draw[->] (input) -- (cond);
  \draw[->] (cond) -- node[right, pos=0.38] {是} (stepa);
  \draw[->] (stepa) -- (stepb);
  \draw[->] (stepb) -- (stepk);

  % k=k+1 后从左侧折返，在输入框与判断框之间重新汇入
  \draw[->] (stepk.west) -- ++(-14mm,0) |- (merge);

  % “否”从菱形右侧绕至输出框上方，再向下进入输出框
  \coordinate[right=16mm of cond] (no-right);
  \coordinate[above=4mm of output] (out-merge);
  \draw (cond.east) -- node[above, pos=0.38] {否} (no-right) |- (out-merge);
  \draw[->] (out-merge) -- (output.north);
  \draw[->] (output) -- (end);
\end{tikzpicture}
